% !TEX root = ../main.tex
% =====================================================================
% ABSTRACT / SOMMARIO
% =====================================================================

\chapter*{Sommario}
\addcontentsline{toc}{chapter}{Sommario}

L'adozione degli strumenti di intelligenza artificiale generativa nei contesti
aziendali è un tema tanto attuale quanto complesso. Da un lato molte voci, 
spesso legate a chi sviluppa o ridistribuisce queste tecnologie, la
presentano come una rivoluzione capace di trasformare il lavoro, ridefinire i
processi e aumentare la produttività su scale fino a poco tempo fa impensabili.
Dall'altro, l'adozione \emph{enterprise} si scontra con una serie di problemi e
implicazioni che non godono ancora della stessa risonanza, riassumibili in tre categorie.

\paragraph{Costi.} Il costo di utilizzo dei modelli LLM commerciali più
performanti è in costante aumento: in caso di adozione poco controllata, la spesa
associata diventa difficile da governare. Allo stesso tempo, stimare il risparmio
reale (idealmente derivante dall'aumento di produttività o dalla riduzione dei costi
operativi) è complesso e tanto più incerto quanto meno
l'implementazione è ``a fuoco'' e puntuale sui processi e casi d'uso.

\paragraph{Etica.} Come noto e dichiarato dai principali sviluppatori, i modelli LLM sono
esposti a bias di varia natura. In contesti aziendali in cui i valori etici sono
parte integrante dell'identità organizzativa, un'adozione indiscriminata rischia
di diffondere nozioni e approcci in aperto contrasto con la condotta
dell'organizzazione.

\paragraph{Sovranità del dato.} Questi strumenti, specialmente se usati in modo
incontrollato o fuori policy dai collaboratori, non offrono garanzie adeguate sulla
confidenzialità e sulla sicurezza dei dati. Questo può mettere a rischio tanto gli utenti
quanto l'organizzazione.

Questa tesi descrive in dettaglio come, nella realtà di Coop Alleanza 3.0, sia
stato implementato un ecosistema di strumenti conversazionali generativi che
affronta e mitiga i problemi appena citati. Descrive inoltre un approccio metodico
e misurabile all'orchestrazione del loro utilizzo, realizzato attraverso un
\emph{LLM Orchestrator} che instrada dinamicamente le richieste dell'utente verso
l'agente specializzato più appropriato. La valutazione sperimentale misura l'accuratezza
dell'orchestratore e la qualità del responso degli agenti verticali (execution accuracy
per il Text-to-SQL, faithfulness per il RAG), documenta le tecniche di \emph{prompt
engineering} applicate nello sviluppo e ottimizzazione degli agenti 
e quantifica il trade-off costo--qualità realizzato dal routing di complessità, f
ornendo infine una stima il più possibile attendibile della sostenibilità 
economica della soluzione.

\clearpage
